% Purpose: Mathematical formulation, algorithmic implementation, and HPC
% orchestration of deep-learning seismic phase pickers (PhaseNet and
% EQTransformer) fine-tuning and regional domain adaptation on OGS
% continuous waveform archives.
% Status: Under Review
% Source-of-truth: OGS/src/ogstrainer.py, OGS/src/ogsconstants.py,
% OGS/utils/Leonardo/Makefile, OGS/utils/Leonardo/LAUNCHME.sh,
% OGS/utils/Leonardo/ktanakah.sh

\providecommand{\bm}[1]{\mathbf{#1}}

\label{chap:training}

Automated seismic monitoring pipelines deployed across dense regional
networks rely fundamentally upon high-sensitivity, robust deep-learning
phase pickers to identify primary ($P$) and secondary ($S$) wave arrival
onsets within continuous multi-channel waveform streams. While deep neural
architectures such as PhaseNet (\cite{PhaseNet}) and Earthquake Transformer
(\ac{EQTransformer}, \cite{EQTransformer}) achieve remarkable pick
precision and recall when evaluated on standard benchmark datasets, their
direct zero-shot deployment in localized seismotectonic domains often
incurs significant performance degradation. This discrepancy stems from
pronounced domain shift between the generalized training distributions and
the specific wave propagation characteristics, crustal velocity structures,
ambient noise environments, and sensor instrumentation of the target
observational domain.

To bridge this operational performance gap, this chapter formulates and
implements an end-to-end, reproducible training and fine-tuning framework
engineered for regional domain adaptation on continuous waveform archives
recorded by the \ac{OGS} \ac{SMINO} network (\cite{SMINO}) across
Northeastern Italy and the Alpine-Dinaric transition zone. The training
system is implemented in the production module \path{OGS/src/ogstrainer.py}
under the unified SeisBench framework (\cite{SeisBench}), bridging
analyst-curated catalog bulletin archives with High-Performance Computing
(\ac{HPC}) infrastructure.

This chapter is structured into seven complementary sections:
\begin{itemize}
  \item Section~\ref{sec:domain_adaptation_theory} establishes the
    theoretical rationale for seismic transfer learning and domain
    adaptation, dissecting the geophysical drivers of covariate shift
    across the complex crustal velocity structures of Northeastern Italy.
  \item Section~\ref{sec:data_extraction_pipeline} details the data
    extraction pipeline, analyst-pick quality filtering based on Hypo71
    weights, automated waveform retrieval, signal conditioning, and
    SeisBench HDF5 container serialization.
  \item Section~\ref{sec:neural_architectures} presents the architectural
    specifications of the PhaseNet 1D U-Net and \ac{EQTransformer}
    hierarchical attention networks alongside their pretrained baseline
    weight repositories.
  \item Section~\ref{sec:data_augmentation} formulates the stochastic data
    augmentation graph, random window selection, and Gaussian probabilistic
    target labeling for arrival time representation.
  \item Section~\ref{sec:loss_and_optimization} derives the multi-channel
    vector cross-entropy loss function, batch-level standard normalization,
    and Adam optimization dynamics.
  \item Section~\ref{sec:checkpointing_and_validation} outlines the
    dataset partitioning strategy, validation protocol, and state-preserving
    checkpoint retention architecture.
  \item Section~\ref{sec:hpc_orchestration} demonstrates the HPC workload
    orchestration on the CINECA Leonardo supercomputer, detailing
    parameterized Makefile rules and SLURM execution wrappers.
\end{itemize}

\paragraph{Evidence status and publication boundary.}
The software modules governing catalog ingestion, waveform conditioning,
stochastic data augmentation, vector cross-entropy loss calculation,
model checkpointing, and SLURM workload generation are demonstrated and
fully verified within \path{OGS/src/ogstrainer.py}, \path{OGS/src/ogsconstants.py},
and \path{OGS/utils/Leonardo/Makefile}. Conversely, the complete multi-epoch
training trajectories, converged regional model weights, and prospective
generalization benchmarks across the multi-year OGS archive represent
active computational campaigns scheduled on Leonardo Booster GPU allocations.
Accordingly, numerical loss trajectories and performance metrics reported
herein denote intended evaluation schemas and architectural baseline
parameters rather than final certified catalog revisions.

% =============================================================================
% SECTION 1: THEORETICAL RATIONALE AND REGIONAL DOMAIN ADAPTATION
% =============================================================================
\section{Theoretical Rationale and Regional Domain Adaptation}
\label{sec:domain_adaptation_theory}

Deep neural network phase pickers have transformed observational
seismology by replacing heuristic, univariate triggering algorithms
\eg{\ac{STA/LTA}, \cite{STALTA}; autoregressive pickers, \cite{LEONARD1999247}}
with multi-channel representation learning. Standard benchmark models are
typically trained on large-scale public repositories such as \ac{SCEDC}
(\cite{SCEDC}) in Southern California, \ac{STEAD} (\cite{STEAD}) globally,
or \ac{INSTANCE} (\cite{INSTANCE}) across the Italian peninsula. As
demonstrated by systematic benchmarking studies (\cite{WhichPicker}), models
trained on one regional dataset often suffer pronounced accuracy degradation,
elevated picking variance, and heightened false-positive rates when evaluated
out-of-domain on unseen seismic networks.

\subsection{Mathematical Formulation of Seismic Domain Shift}
\label{subsec:domain_shift_math}

In statistical machine learning, a seismic phase picking domain
$\mathcal{D} = \{\mathcal{X}, P(\mathbf{X})\}$ is defined by a feature space
$\mathcal{X} \subset \mathbb{R}^{C \times L}$ of multi-component seismic
waveform time series and a marginal probability distribution $P(\mathbf{X})$,
where $C = 3$ denotes the three orthogonal recording components (vertical $Z$,
north-south $N$, east-west $E$) and $L$ denotes the window sequence length.
A phase picking task $\mathcal{T} = \{\mathcal{Y}, P(\mathbf{Y} \mid \mathbf{X})\}$
consists of a label space $\mathcal{Y} \subset [0, 1]^{K \times L}$ representing
continuous temporal probability trajectories for $K = 3$ target classes
$\{P, S, \mathrm{Noise}\}$ and an objective conditional probability distribution
$P(\mathbf{Y} \mid \mathbf{X})$ modeled by a neural network
$f_{\bm{\Theta}}: \mathcal{X} \to \mathcal{Y}$ parameterized by weights $\bm{\Theta}$.

During supervised pretraining, the neural network optimizes its parameters on a
source domain $\mathcal{D}_s = \{\mathcal{X}_s, P_s(\mathbf{X})\}$ under source
task $\mathcal{T}_s$ by minimizing an empirical risk objective:
\begin{equation}
  \bm{\Theta}^* = \argmin_{\bm{\Theta}} \mathbb{E}_{(\mathbf{X}, \mathbf{Y}) \sim \mathcal{D}_s}
  \left[ \mathcal{L}\bigl( f_{\bm{\Theta}}(\mathbf{X}), \mathbf{Y} \bigr) \right].
  \label{eq:source_risk_minimization}
\end{equation}

When the trained model $f_{\bm{\Theta}^*}$ is deployed on a regional target
domain $\mathcal{D}_t = \{\mathcal{X}_t, P_t(\mathbf{X})\}$, such as the OGS
network in Northeastern Italy, domain shift manifests across two primary
statistical regimes:
\begin{enumerate}
  \item \textbf{Covariate Shift} ($P_s(\mathbf{X}) \neq P_t(\mathbf{X})$ while
    $P_s(\mathbf{Y} \mid \mathbf{X}) \approx P_t(\mathbf{Y} \mid \mathbf{X})$):
    The marginal distribution of seismic wavefields differs between domains due
    to differing ambient seismic noise spectra, site amplification, and sensor
    transfer functions, while the physical relationship between wave arrival
    onsets and particle motion remains invariant.
  \item \textbf{Concept Drift / Conditional Shift} ($P_s(\mathbf{Y} \mid \mathbf{X}) \neq P_t(\mathbf{Y} \mid \mathbf{X})$):
    The conditional distribution of phase arrival probabilities diverges because
    differing regional crustal velocity gradients, structural discontinuities,
    and scattering regimes distort waveform signatures, generating emergent phase
    arrivals, complex coda conversions, and multi-path reflections not captured in
    the source training distribution.
\end{enumerate}

The discrepancy between source and target feature representations can be
formalized via the Kullback--Leibler (\ac{KL}) divergence between the induced
marginal feature distributions:
\begin{equation}
  D_{\mathrm{KL}}\bigl( P_s(\mathbf{X}) \,\|\, P_t(\mathbf{X}) \bigr) =
  \int_{\mathcal{X}} P_s(\mathbf{X}) \log\left( \frac{P_s(\mathbf{X})}{P_t(\mathbf{X})} \right) \mathrm{d}\mathbf{X}.
  \label{eq:kl_divergence}
\end{equation}
Whenever $D_{\mathrm{KL}}(P_s \,\|\, P_t) > 0$, the theoretical generalization
bound on target empirical error $\mathcal{E}_t(f_{\bm{\Theta}})$ increases
proportionally to the distribution divergence, leading to missed weak onsets
and spurious picks on complex coda.

\subsection{Geophysical Drivers of Domain Shift in Northeastern Italy}
\label{subsec:geophysical_drivers}

The target operational environment of the \ac{OGS} \ac{SMINO} network encloses
an active tectonic collision zone spanning the Eastern Southern Alps and the
External Dinarides. Several regional geophysical factors induce pronounced
domain discrepancy relative to standard training corpora:

\begin{itemize}
  \item \textbf{Complex Crustal Velocity Discontinuities}: The crust in
    Northeastern Italy exhibits extreme lateral and vertical heterogeneity.
    The Moho depth deepens rapidly from approximately $30\,\mathrm{km}$ beneath
    the Venetian-Friulian alluvial plain to over $45\,\mathrm{km}$ beneath the
    Alpine metamorphic root (\cite{Bressan1998}). Thrust-and-fold belts in the
    Carnic and Julian Prealps create sharp acoustic impedance contrasts,
    inducing complex refracted head waves ($P_n$, $S_n$) and post-critical
    Moho reflections ($P_mP$, $S_mS$) with emergent onsets that confound models
    trained on simpler tabular crustal structures.
  \item \textbf{High-Frequency Attenuation and Scattering}: The fractured
    carbonate platforms and karst lithologies characterizing the Friuli region
    exhibit high intrinsic and scattering attenuation. The seismic quality factor
    follows a frequency-dependent power law $Q(f) = Q_0 f^\alpha$, with regional
    estimates indicating significant high-frequency damping ($f > 15\,\mathrm{Hz}$)
    and severe near-surface high-frequency spectral decay governed by the
    Anderson-Hough attenuation parameter $\kappa_0$. High-frequency $P$-wave
    energy is rapidly dispersed into coda waves, softening impulse onsets into
    gradual ramps.
  \item \textbf{Topographic and Valley Site Amplification}: Many \ac{OGS} stations
    are sited in narrow Alpine valleys filled with soft Quaternary alluvium or on
    steep topography. Trapped sedimentary valley waves and topographic ridge
    scattering create monochromatic resonance bands and energetic S-wave coda
    reverberations that mimic genuine microseismic arrivals.
  \item \textbf{Heterogeneous Multi-Network Instrumentation}: The OGS monitoring
    archive ingests telemetry from 18 distinct regional and international networks,
    operating high-sensitivity broadband velocimeters (Streckeisen STS-2, Nanometrics
    Trillium 120/240) alongside strong-motion accelerometers (Kinemetrics Episensor)
    with disparate sensor poles, zeroes, dynamic ranges, and digitizer sampling
    frequencies ($50$, $100$, and $200\,\mathrm{Hz}$).
\end{itemize}

\subsection{Inductive Transfer Learning vs Fine-Tuning Strategies}
\label{subsec:transfer_strategies}

Addressing domain discrepancy requires selecting an appropriate transfer learning
strategy. In seismology, two principal paradigms exist:
\begin{enumerate}
  \item \textbf{Full Retraining from Random Initialization}: Training deep
    neural networks from scratch requires massive ground-truth training corpora
    ($N \ge 10^5\text{--}10^6$ labeled traces). Regional seismic bulletins,
    including the \ac{OGS} catalog (\cite{MagrinEtAl2026RevisedOGSCatalog}),
    possess high-quality analyst labels but with smaller total sample volumes
    ($N \approx 10^3\text{--}10^4$ events). Training complex architectures from
    scratch on regional corpora inevitably leads to catastrophic overfitting,
    where the network memorizes station-specific noise signatures rather than
    learning generalized wavefield kinematics.
  \item \textbf{Inductive Transfer Learning via Parameter Fine-Tuning}:
    Under this approach, the network weights $\bm{\Theta}$ are warm-started
    from a model pretrained on large benchmark corpora ($\bm{\Theta}^{(0)} \leftarrow \bm{\Theta}_{\mathrm{pretrained}}$).
    The early convolutional layers, which extract universal low-level waveform
    features such as onset slopes, frequency content, and amplitude envelopes,
    remain largely preserved. Fine-tuning with a reduced learning rate
    $\eta \in [10^{-3}, 10^{-2}]$ adapts the intermediate representations and
    deeper decision layers to regional attenuation, site resonance, and complex
    travel-time trajectories.
\end{enumerate}

The framework implemented in \path{OGS/src/ogstrainer.py} adopts inductive
transfer learning, enabling systematic fine-tuning from multiple pretrained
baselines to identify the optimal regional domain adaptation trajectory.

% =============================================================================
% SECTION 2: DATA EXTRACTION AND STREAM CONDITIONING PIPELINE
% =============================================================================
\section{Data Extraction and Stream Conditioning Pipeline}
\label{sec:data_extraction_pipeline}

The training dataset generation engine bridges the authoritative, analyst-curated
OGS seismic catalog with the continuous multi-channel waveform repository. The
workflow is coordinated by the \texttt{OGSTrainer.build\_dataset()} method in
\path{OGS/src/ogstrainer.py}.

\subsection{Analyst Label Ingestion and Pick Quality Filtering}
\label{subsec:pick_filtering}

Ground-truth phase labels are retrieved from the object-oriented \texttt{OGSCatalog}
interface across a user-specified Gregorian date window $[t_{\mathrm{start}}, t_{\mathrm{end}}]$.
For each calendar day within the target observation window, pick records and event
hypocenters are loaded into memory and grouped by unique pick index (\texttt{index})
and station identifier (\texttt{station}).

To ensure that supervised gradient updates are driven solely by high-precision,
unambiguous seismic onsets, an analyst pick quality filter is enforced. The
\ac{OGS} bulletin encodes arrival-time uncertainties using the standard Hypo71
(\cite{Hypo71}) integer weighting scheme $w \in \{0, 1, 2, 3, 4, 5\}$, defined
in \path{OGS/src/ogsconstants.py}:
\begin{equation}
  \sigma_t(w) =
  \begin{cases}
    0.01\,\unit{\second}, & w = 0 \quad (\text{Impulsive, very sharp onset}), \\
    0.04\,\unit{\second}, & w = 1 \quad (\text{Clear, distinct onset}), \\
    0.20\,\unit{\second}, & w = 2 \quad (\text{Fairly clear onset}), \\
    1.00\,\unit{\second}, & w = 3 \quad (\text{Emergent onset}), \\
    5.00\,\unit{\second}, & w = 4 \quad (\text{Poor quality, highly uncertain}), \\
    25.0\,\unit{\second}, & w = 5 \quad (\text{Unusable / discarded phase}).
  \end{cases}
  \label{eq:hypo71_weights}
\end{equation}

In \texttt{ogstrainer.py}, picks are strictly filtered to retain only high-confidence
arrivals with weights $w \le 2$:
\begin{lstlisting}[language=python,basicstyle=\scriptsize\ttfamily]
station_picks = station_picks[station_picks[OGS_C.PHASE_STR].notna()]
station_picks = station_picks[
    station_picks[OGS_C.WEIGHT_STR].astype(float) <= 2
]
if len(station_picks) <= 1:
    continue
\end{lstlisting}
This filtering step eliminates emergent and questionable picks ($w \ge 3$),
guaranteeing that the temporal target distributions represent physical arrivals
with timing precision $\sigma_t \le 0.20\,\unit{\second}$. Furthermore, stations
recording only an isolated pick without a complementary phase arrival are
pruned from training pair generation to enforce balanced multi-phase representation.

\subsection{Waveform Stream Windowing and Automated Retrieval}
\label{subsec:waveform_windowing}

For each validated pick, the corresponding continuous waveform files are located
within the hierarchical miniSEED archive structured by year, month, and day:
\begin{equation}
  \path{<WAVE_PATH>/<YYYY>/<MM>/<DD>/<NET>.<STA>.*.*__<START_TIME>__<END_TIME>.mseed}
  \label{eq:mseed_path_pattern}
\end{equation}

The waveform segment is extracted with an operational temporal padding of
$\pm 60\,\unit{\second}$ centered around the pick timestamp $t_{\mathrm{pick}}$,
defined by the time delta constant \texttt{PICK\_TRAIN\_OFFSET = td(seconds=60)}
in \path{OGS/src/ogsconstants.py}:
\begin{equation}
  \mathcal{W}_{\mathrm{extract}} = \left[ t_{\mathrm{pick}} - 60\,\unit{\second}, \;
  t_{\mathrm{pick}} + 60\,\unit{\second} \right].
  \label{eq:extraction_window}
\end{equation}
This $120\,\unit{\second}$ temporal window guarantees that both primary and secondary
arrivals, the pre-P noise baseline, and the complete decaying coda are contained
within the raw trace prior to subsequent stochastic cropping.

When continuous waveforms are absent from local scratch mounts, \texttt{OGSTrainer}
provides automated fallback retrieval via \texttt{\_download\_waveform()}, which
executes \path{OGS/src/ogsdownloader.py} as an asynchronous subprocess querying
distributed \ac{FDSN} data nodes (\cite{SMINO}) across INGV, SED, ARSO, and OGS
repositories.

\subsection{Signal Conditioning Pipeline}
\label{subsec:signal_conditioning}

Raw waveforms recorded by heterogeneous digitizers exhibit DC offsets, linear
instrumental drift, and high-frequency anthropogenic contamination. Prior to
dataset compilation, every extracted trace is conditioned via the deterministic
preprocessing method \texttt{get\_clean_trace()}:
\begin{equation}
  x(t) \xrightarrow{\text{Detrend}} x_1(t) \xrightarrow{\text{Demean}} x_2(t)
  \xrightarrow{\text{Taper}} x_3(t) \xrightarrow{\text{Bandpass}} x_4(t)
  \xrightarrow{\text{Resample}} x_5(t).
  \label{eq:preprocessing_chain}
\end{equation}

The mathematical operations comprising this pipeline include:
\begin{enumerate}
  \item \textbf{Linear Detrending}: Fitting and subtracting an optimal first-order
    polynomial $p(t) = m t + c$ via least-squares regression:
    \begin{equation}
      x_1(t) = x(t) - (m t + c), \quad \text{where} \quad (m, c) = \argmin_{m', c'}
      \sum_{k=1}^N \bigl( x(t_k) - (m' t_k + c') \bigr)^2.
    \end{equation}
  \item \textbf{Mean Subtraction (Demeaning)}: Centering the trace to zero mean:
    \begin{equation}
      x_2(t) = x_1(t) - \frac{1}{N} \sum_{k=1}^N x_1(t_k).
    \end{equation}
  \item \textbf{Cosine Tapering}: Applying a symmetric cosine bell taper
    (Tukey window with $5\%$ edge taper, $\alpha = 0.05$) to smoothly attenuate
    trace boundaries to zero, preventing Gibbs oscillations and spectral leakage
    during subsequent Fourier domain filtering:
    \begin{equation}
      w(t) =
      \begin{cases}
        \frac{1}{2}\left[ 1 - \cos\left( \frac{\pi t}{\alpha T} \right) \right], & 0 \le t < \alpha T, \\
        1, & \alpha T \le t \le (1 - \alpha) T, \\
        \frac{1}{2}\left[ 1 - \cos\left( \frac{\pi (T - t)}{\alpha T} \right) \right], & (1 - \alpha) T < t \le T.
      \end{cases}
    \end{equation}
  \item \textbf{4th-Order Zero-Phase Butterworth Bandpass Filtering}: Filtering
    between $f_{\mathrm{min}} = 0.1\,\unit{\hertz}$ and $f_{\mathrm{max}} = 20.0\,\unit{\hertz}$.
    The 4th-order zero-phase implementation applies forward-backward filtering
    to achieve zero group-delay distortion, preserving exact phase onset timing:
    \begin{equation}
      |H(f)|^2 = \frac{1}{1 + \left( \frac{f^2 - f_0^2}{f \cdot \Delta f} \right)^{2 n}}, \quad n = 4,
    \end{equation}
    where $\Delta f = f_{\mathrm{max}} - f_{\mathrm{min}} = 19.9\,\unit{\hertz}$
    and $f_0 = \sqrt{f_{\mathrm{min}} f_{\mathrm{max}}}$.
  \item \textbf{Anti-Aliasing Resampling}: Resampling the conditioned stream to a
    normalized sampling frequency $f_s = 100.0\,\unit{\hertz}$ ($\Delta t = 0.01\,\unit{\second}$)
    using polyphase rational interpolation with an integrated anti-aliasing low-pass
    filter at the Nyquist cutoff ($f_{\mathrm{Nyq}} = 50.0\,\unit{\hertz}$).
\end{enumerate}

\subsection{SeisBench Waveform Dataset Serialization}
\label{subsec:seisbench_serialization}

Conditioned traces and matched catalog records are compiled into a flat tabular
index, \path{training_dataset.csv}, storing comprehensive trace metadata:
source identifiers, network, station, channel, event hypocentral coordinates
(latitude, longitude, depth), origin time, local magnitude ($M_{\mathrm{L}}$),
phase type ($P$ or $S$), arrival timestamp, peak trace amplitude, and Hypo71
weight.

To support high-throughput, random-access mini-batch loading across parallel
GPU workers, this flat dataset is converted into SeisBench standardized format
using \texttt{seisbench.data.WaveformDataWriter}. The writer generates:
\begin{enumerate}
  \item \path{waveforms.hdf5}: An optimized hierarchical HDF5 container storing
    multi-component waveform arrays under chunked, contiguous binary datasets
    with blosc compression, enabling zero-copy slicing of arbitrary window offsets.
  \item \path{metadata.csv}: A standardized columnar schema indexing trace
    indices, sample arrival positions (\eg{\texttt{trace\_p\_arrival\_sample}} and
    \texttt{trace\_s\_arrival\_sample}), epicentral distances, and source magnitudes.
\end{enumerate}
Once serialized, the dataset is ingested into \texttt{seisbench.data.WaveformDataset},
providing a high-performance abstraction that decouples raw waveform storage from
neural pipeline execution.

% =============================================================================
% SECTION 3: NEURAL MODEL ARCHITECTURES AND PRETRAINED BASELINES
% =============================================================================
\section{Neural Model Architectures and Pretrained Baselines}
\label{sec:neural_architectures}

The training engine interfaces with two distinct neural network topologies
registered in the \texttt{MODEL\_REGISTRY} of \path{OGS/src/ogstrainer.py}:
PhaseNet (\cite{PhaseNet}) and Earthquake Transformer (\ac{EQTransformer},
\cite{EQTransformer}).

\subsection{PhaseNet: 1D U-Net with Lateral Skip Connections}
\label{subsec:phasenet_arch}

PhaseNet adapts the U-Net architecture from biomedical image segmentation to
one-dimensional multi-channel time series. The network consists of a contracting
path (encoder) and an expanding path (decoder) joined by lateral skip connections,
as formulated in Chapter~\ref{chap:pipeline}.

\begin{itemize}
  \item \textbf{Input Dimensionality}: Three-component seismic waveform
    $\mathbf{X} \in \mathbb{R}^{3 \times L}$, where $L = 3{,}001$ samples
    ($30.01\,\unit{\second}$ at $100\,\unit{\hertz}$).
  \item \textbf{Contracting Encoder}: Four sequential downsampling stages.
    Each stage comprises a 1D convolution with a large kernel size $k = 7$,
    stride $s = 1$, batch normalization, and a rectified linear unit (ReLU)
    activation function, followed by a strided 1D convolution with kernel $k=7$
    and stride $s = 2$ that halves the temporal resolution while doubling the
    feature channel depth:
    \begin{equation}
      C_l = 8 \times 2^l, \quad L_l = \left\lfloor \frac{L_{l-1} + 1}{2} \right\rfloor, \quad l \in \{1, 2, 3, 4\}.
    \end{equation}
  \item \textbf{Bottleneck}: A deep contextual representation layer with 128
    feature channels operating on the downsampled feature map ($L_4 \approx 188$
    samples), capturing long-period waveform envelopes and broad temporal context.
  \item \textbf{Expanding Decoder}: Four symmetric upsampling stages. Each
    stage employs transposed 1D convolutions with stride $s = 2$ to double the
    temporal dimension, concatenates the resulting feature map laterally with the
    corresponding high-resolution representation from the contracting encoder
    via skip connections, and processes the merged representation through 1D
    convolutions and ReLU activations.
  \item \textbf{Output Layer}: A $1 \times 1$ pointwise convolution mapping
    the 8-channel reconstructed feature map to three logits, followed by a
    temporal softmax projection producing posterior probabilities:
    \begin{equation}
      \hat{y}_{c}(t) = \frac{\exp\bigl( z_c(t) \bigr)}{\sum_{k \in \{P, S, N\}} \exp\bigl( z_k(t) \bigr)}, \quad c \in \{P, S, N\},
      \label{eq:phasenet_softmax}
    \end{equation}
    where $\hat{y}_P(t)$ is the P-wave probability, $\hat{y}_S(t)$ is the S-wave
    probability, and $\hat{y}_N(t)$ is the background noise probability.
\end{itemize}

The combination of deep downsampling and lateral skip connections enables
PhaseNet to resolve large temporal receptive fields ($\approx 10\text{--}15\,\unit{\second}$)
while preserving high-frequency onset timing down to the individual sample level.

\subsection{EQTransformer: Attentive Deep Hierarchical Architecture}
\label{subsec:eqt_arch}

Earthquake Transformer (\ac{EQTransformer}, \cite{EQTransformer}) combines
1D convolutional residual networks, bidirectional Long Short-Term Memory
(\ac{LSTM}) recurrent networks, and multi-head self-attention mechanisms to
jointly perform earthquake detection and phase picking.

\begin{itemize}
  \item \textbf{Input Dimensionality}: Three-component waveform $\mathbf{X} \in \mathbb{R}^{3 \times L}$
    with sequence length $L = 6{,}000$ samples ($60.00\,\unit{\second}$ at $100\,\unit{\hertz}$).
  \item \textbf{1D ResNet Feature Extractor}: A hierarchy of 1D residual blocks
    operating with strided convolutions that downsamples the input sequence by a
    cumulative factor of 8 to $L_{\mathrm{down}} = 750$ time steps, extracting
    local temporal features while mitigating vanishing gradient degradation.
  \item \textbf{Bidirectional LSTM Sequence Modeling}: A bidirectional LSTM layer
    processes the compressed temporal features in forward and backward directions,
    capturing long-range contextual dependencies across the waveform sequence:
    \begin{equation}
      \mathbf{h}_t = \left[ \overrightarrow{\mathrm{LSTM}}(\mathbf{x}_t, \overrightarrow{\mathbf{h}}_{t-1}) \,\|\,
      \overleftarrow{\mathrm{LSTM}}(\mathbf{x}_t, \overleftarrow{\mathbf{h}}_{t+1}) \right].
    \end{equation}
  \item \textbf{Multi-Head Self-Attention}: Scaled dot-product attention
    mechanisms allow the network to dynamically assign attention weights across
    distant temporal segments:
    \begin{equation}
      \mathrm{Attention}(\mathbf{Q}, \mathbf{K}, \mathbf{V}) =
      \mathrm{softmax}\left( \frac{\mathbf{Q}\mathbf{K}^T}{\sqrt{d_k}} \right)\mathbf{V},
    \end{equation}
    where query $\mathbf{Q}$, key $\mathbf{K}$, and value $\mathbf{V}$ projections
    enable the model to correlate initial P-wave particle motion onsets with subsequent
    S-wave shear arrivals and coda energy distributions.
  \item \textbf{Multi-Task Decoders}: Three independent feed-forward decoder heads
    with separate attention modules output continuous probability trajectories for:
    (i) earthquake detection ($\hat{y}_{\mathrm{det}}$), (ii) P-phase arrival
    ($\hat{y}_P$), and (iii) S-phase arrival ($\hat{y}_S$).
\end{itemize}

\subsection{Pretrained Baseline Model Repositories}
\label{subsec:pretrained_repositories}

The \texttt{ogstrainer.py} CLI exposes five benchmark pretrained weight variants
via the \texttt{-s / --dataset} flag, corresponding to major public seismic
corpora distributed within the SeisBench repository:
\begin{enumerate}
  \item \texttt{instance} (\cite{INSTANCE}): Pretrained on the Italian Seismic
    Dataset for Machine Learning, comprising $1.15 \times 10^6$ three-component
    traces recorded across Italy. Because it encompasses the tectonic regime of
    peninsular Italy, it serves as the primary regional baseline for transfer learning.
  \item \texttt{stead} (\cite{STEAD}): Pretrained on the Stanford Earthquake
    Dataset, containing $1.2 \times 10^6$ global earthquake and noise traces with
    broad geometric diversity.
  \item \texttt{scedc} (\cite{SCEDC}): Pretrained on Southern California Earthquake
    Data Center records, dominated by strike-slip faulting mechanisms and shallow
    crustal structures.
  \item \texttt{original} (\cite{PhaseNet,EQTransformer}): Initial author weights
    trained on Northern California (NCEDC) data for PhaseNet and mixed global/California
    data for EQTransformer.
  \item \texttt{adriaarray}: Pretrained on regional seismic waveforms collected
    during the recent AdriaArray scientific initiative across the Adriatic-Pannonian
    collision zone.
\end{enumerate}

Table~\ref{tab:pretrained_corpora_comparison} synthesizes the key seismotectonic
and instrumental properties of these five baseline training corpora.

\begin{table}[htbp]
  \centering
  \small
  \caption{Comparative synthesis of benchmark training corpora and pretrained
    baseline weight repositories available for regional transfer learning.}
  \label{tab:pretrained_corpora_comparison}
  \begin{threeparttable}
    \begin{tabularx}{\linewidth}{l l l c >{\raggedright\arraybackslash}X}
      \toprule
      \textbf{Corpus Key} & \textbf{Geographic Scope} & \textbf{Tectonic Regime} & \textbf{Traces [$10^3$]} & \textbf{Transfer Learning Compatibility to NE Italy} \\
      \midrule
      \texttt{instance}   & Italy (INGV)              & Apennines / Alps         & $1{,}150$                 & High; shares Italian network instrumentation and regional wave propagation terms. \\
      \texttt{adriaarray} & Central/SE Europe         & Alpine-Dinaric-Pannonian & $\sim 450$                & High; overlaps target regional tectonic boundaries and seismic stations. \\
      \texttt{stead}      & Global                    & Diverse / Global Mix     & $1{,}200$                 & Moderate; excellent SNR diversity but lacks regional crustal wavefield specifics. \\
      \texttt{scedc}      & Southern California       & Strike-slip (San Andreas)& $750$                     & Moderate; distinct sediment basin propagation and shallow seismogenic depths. \\
      \texttt{original}   & Northern California       & Strike-slip / Subduction & $\sim 700$                & Moderate to Low; baseline author reference with California-specific noise textures. \\
      \bottomrule
    \end{tabularx}
    \begin{tablenotes}[flushleft]
      \footnotesize
      \item Target \& Scope: Public seismological training corpora available within
        SeisBench (\cite{SeisBench}) and integrated into \path{OGS/src/ogstrainer.py}.
      \item Notice that \texttt{instance} and \texttt{adriaarray} provide the closest
        crustal velocity and instrumental alignment with the target Alpine-Dinaric domain.
      \item Evidence Anchor: Configured via \path{OGS/src/ogsconstants.py} constants
        \texttt{INSTANCE\_STR}, \texttt{ADRIAARRAY\_STR}, \texttt{STEAD\_STR},
        \texttt{SCEDC\_STR}, and \texttt{ORIGINAL\_STR}.
    \end{tablenotes}
  \end{threeparttable}
\end{table}

% =============================================================================
% SECTION 4: DATA AUGMENTATION AND PROBABILISTIC PHASE LABELING
% =============================================================================
\section{Data Augmentation and Probabilistic Phase Labeling}
\label{sec:data_augmentation}

To prevent the deep neural networks from memorizing absolute phase positions or
overfitting to localized station noise textures, \path{OGS/src/ogstrainer.py}
implements an advanced stochastic data augmentation and target generation graph
via \texttt{seisbench.generate}.

\subsection{Stochastic Augmentation Graph via SeisBench Generate}
\label{subsec:augmentation_graph}

The augmentation pipeline is constructed dynamically in \texttt{\_build\_augmentations()}
and integrated into \texttt{seisbench.generate.GenericGenerator}. The operational
sequence comprises five sequential stochastic transformations:

\begin{enumerate}
  \item \textbf{WindowAroundSample}: Selects a wide candidate window guaranteed to
    enclose a verified catalog arrival. Rather than centering the pick deterministically,
    the pick is randomly positioned within a window of length $2 \times L_{\mathrm{in}}$
    with $L_{\mathrm{in}}$ samples preceding the pick, using a variable strategy:
    \begin{lstlisting}[language=python,basicstyle=\scriptsize\ttfamily]
sbg.WindowAroundSample(
    list(PHASE_DICT.keys()),
    samples_before=windowlen,
    windowlen=windowlen * 2,
    selection="random",
    strategy="variable",
)
    \end{lstlisting}
    This prevents the neural network from learning a trivial spatial prior
    centered at the middle of the window.
  \item \textbf{RandomWindow}: Performs a secondary stochastic crop to extract
    the exact model input sequence length ($L=3{,}001$ for PhaseNet, $L=6{,}000$
    for EQTransformer) with edge zero-padding (\texttt{strategy="pad"}):
    \begin{lstlisting}[language=python,basicstyle=\scriptsize\ttfamily]
sbg.RandomWindow(windowlen=windowlen, strategy="pad")
    \end{lstlisting}
    This exposes the model to arbitrary arrival offsets, including picks situated
    near the start or end of the inference window.
  \item \textbf{ChangeDtype (Input $\mathbf{X}$)}: Explicitly casts waveform sample
    tensors to 32-bit floating point precision (\texttt{np.float32}), matching
    PyTorch accelerator requirements and preventing double-precision memory overhead.
  \item \textbf{ProbabilisticLabeller}: Generates smooth Gaussian probability
    trajectories centered at the exact fractional sample positions of the P and S
    arrivals, as detailed in Section~\ref{subsec:probabilistic_labeling}.
  \item \textbf{ChangeDtype (Target $\mathbf{Y}$)}: Casts the output label probability
    tensor $\mathbf{Y}$ to \texttt{np.float32} precision.
\end{enumerate}

\subsection{Probabilistic Phase Label Formulation}
\label{subsec:probabilistic_labeling}

A fundamental challenge in training neural phase pickers lies in target arrival
representation. Formulating phase picking as discrete Dirac impulse classification
($y_c(t) = 1$ if $t = t_{\mathrm{arr}}$, $y_c(t) = 0$ otherwise) incurs two severe
mathematical deficiencies:
\begin{enumerate}
  \item \textbf{Extreme Class Imbalance}: Within an input window of $L = 3{,}001$
    samples, a single discrete impulse corresponds to a positive class prior
    $P(y=1) \approx 3.3 \times 10^{-4}$. Gradients generated by cross-entropy loss
    are overwhelmingly dominated by background noise samples, causing optimization
    instability.
  \item \textbf{Absence of Gradient Distance Feedback}: A Dirac target treats a
    prediction error of 1 sample ($\qty{10}{\milli\second}$) identically to an
    error of 1,000 samples ($\qty{10}{\second}$), assigning zero probability
    density and zero smooth gradient outside the single arrival sample.
\end{enumerate}

To establish a smooth convex optimization manifold that mirrors physical picking
uncertainty, \texttt{ProbabilisticLabeller} parameterizes phase targets as continuous
univariate Gaussian bell curves centered at arrival sample $t_{\mathrm{arr}}$:
\begin{equation}
  y_c(t) = \exp\left( -\frac{(t - t_{\mathrm{arr}})^2}{2\sigma^2} \right), \quad c \in \{P, S\},
  \label{eq:gaussian_target_bell}
\end{equation}
where the standard deviation $\sigma$ is set to $\sigma = 30$ samples ($0.30\,\unit{\second}$
at $f_s = 100\,\unit{\hertz}$). The Full Width at Half Maximum (\ac{FWHM}) of the
resulting target distribution is:
\begin{equation}
  \mathrm{FWHM} = 2\sqrt{2\ln 2}\,\sigma \approx 2.3548 \times 30\,\text{samples} \approx 70.64\,\text{samples} \approx 0.706\,\unit{\second}.
  \label{eq:fwhm_calculation}
\end{equation}
This $0.71\,\unit{\second}$ effective probability envelope comfortably spans the
$0.20\,\unit{\second}$ empirical timing uncertainty threshold associated with
Hypo71 quality weight $w \le 2$ (Section~\ref{subsec:pick_filtering}), providing
the network with smooth, non-vanishing gradient guidance proportional to temporal
offset. Target values are truncated to zero when $y_c(t) < 10^{-4}$.

\subsection{Noise Channel Normalization and Multi-Label Consistency}
\label{subsec:noise_channel_norm}

For three-class networks (PhaseNet), the background noise channel probability
$y_N(t)$ is defined as the complementary upper envelope of the seismic phase
probabilities:
\begin{equation}
  y_N(t) = 1 - \max\bigl( y_P(t), y_S(t) \bigr).
  \label{eq:noise_channel_complement}
\end{equation}

By construction, this formulation enforces two vital mathematical properties:
\begin{equation}
  \max_{c \in \{P, S, N\}} y_c(t) = 1, \quad \text{and} \quad \sum_{c \in \{P, S, N\}} y_c(t) \ge 1, \quad \forall t \in [1, L].
  \label{eq:label_convexity}
\end{equation}
Far from arrival times, $y_P(t) = 0$ and $y_S(t) = 0$, driving $y_N(t) = 1.0$.
In the immediate vicinity of an arrival, $y_c(t_{\mathrm{arr}}) = 1.0$, which drives
the noise target $y_N(t_{\mathrm{arr}}) = 0.0$. In transition regimes between
adjacent arrivals ($S - P$ intervals), the formulation ensures smooth probability
partitioning without artificial discontinuity.

% =============================================================================
% SECTION 5: MATHEMATICAL LOSS FORMULATION AND OPTIMIZATION DYNAMICS
% =============================================================================
\section{Mathematical Loss Formulation and Optimization Dynamics}
\label{sec:loss_and_optimization}

Neural network optimization is governed by the vector cross-entropy loss function
implemented in \texttt{loss\_fn()} of \path{OGS/src/ogstrainer.py}, integrated
with dynamic GPU waveform normalization and deterministic gradient descent.

\subsection{Multi-Channel Vector Cross-Entropy Loss Derivation}
\label{subsec:loss_derivation}

Let $\mathbf{Y} \in [0, 1]^{B \times K \times L}$ denote the ground-truth probabilistic
target tensor and $\hat{\mathbf{Y}} \in [0, 1]^{B \times K \times L}$ denote the
model-predicted posterior probability tensor for a mini-batch of $B$ training
examples, $K = 3$ output classes ($c \in \{P, S, N\}$), and sequence length $L$.
The predicted posterior probabilities are generated by the final temporal softmax
or sigmoid projection layer.

The vector cross-entropy loss per time sample $t \in \{1, \dots, L\}$ and class
$c \in \{1, \dots, K\}$ for batch element $b \in \{1, \dots, B\}$ is formulated as:
\begin{equation}
  h_{b,c,t} = y_{b,c,t} \log\bigl( \hat{y}_{b,c,t} + \epsilon \bigr),
  \label{eq:sample_entropy}
\end{equation}
where $\epsilon = 10^{-5}$ is a strictly positive numerical regularization constant
preventing mathematical divergence ($\log 0 \to -\infty$) when model predictions
approach zero.

The dimensional aggregation follows the SeisBench convention:
\begin{enumerate}
  \item \textbf{Mean over Sequence Length ($L$)}: The cross-entropy is averaged
    across the temporal sample dimension (axis $-1$) to make the loss invariant to
    arbitrary input window lengths ($L=3{,}001$ versus $L=6{,}000$):
    \begin{equation}
      \bar{h}_{b,c} = \frac{1}{L} \sum_{t=1}^{L} y_{b,c,t} \log\bigl( \hat{y}_{b,c,t} + \epsilon \bigr).
    \end{equation}
  \item \textbf{Sum over Class Channels ($K$)}: The loss is summed across the $K$
    mutually dependent target classes (axis $-1$ after sample reduction):
    \begin{equation}
      H_b = \sum_{c \in \{P, S, N\}} \bar{h}_{b,c} = \sum_{c \in \{P, S, N\}}
      \frac{1}{L} \sum_{t=1}^{L} y_{b,c,t} \log\bigl( \hat{y}_{b,c,t} + \epsilon \bigr).
    \end{equation}
  \item \textbf{Mean over Mini-Batch ($B$)}: The aggregated entropy is averaged
    across all $B$ independent samples within the mini-batch:
    \begin{equation}
      \mathcal{L}(\mathbf{Y}, \hat{\mathbf{Y}}) = - \frac{1}{B} \sum_{b=1}^{B} H_b =
      - \frac{1}{B} \sum_{b=1}^{B} \sum_{c \in \{P, S, N\}} \frac{1}{L} \sum_{t=1}^{L}
      y_{b,c,t} \log\bigl( \hat{y}_{b,c,t} + \epsilon \bigr).
      \label{eq:final_vector_loss}
    \end{equation}
\end{enumerate}

This formulation maps directly to the vectorized implementation in \path{OGS/src/ogstrainer.py}:
\begin{lstlisting}[language=python,basicstyle=\scriptsize\ttfamily]
def loss_fn(y_pred, y_true, eps=1e-5):
    h = y_true * torch.log(y_pred + eps)
    h = h.mean(-1).sum(-1)  # Mean along sample dim, sum along pick dim
    h = h.mean()            # Mean over batch axis
    return -h
\end{lstlisting}

The analytical partial derivative with respect to predicted logit $\hat{y}_{b,c,t}$
evaluates to:
\begin{equation}
  \frac{\partial \mathcal{L}}{\partial \hat{y}_{b,c,t}} =
  - \frac{1}{B L} \cdot \frac{y_{b,c,t}}{\hat{y}_{b,c,t} + \epsilon}.
  \label{eq:loss_gradient}
\end{equation}
Because the denominator is bounded from below by $\epsilon = 10^{-5}$, gradient
magnitudes remain strictly bounded throughout all optimization stages, preventing
exploding gradients during the initial fine-tuning epochs.

\subsection{Batch-Level Standard Waveform Normalization}
\label{subsec:batch_normalization}

Seismic waveform amplitudes vary across several orders of magnitude due to
earthquake magnitude scaling ($\Delta M_{\mathrm{L}} = 1$ corresponds to a factor
of 10 in ground amplitude) and geometric spreading attenuation. To prevent large-magnitude
events from dominating gradient updates, dynamic trace standard normalization is
executed on the GPU device via \texttt{model.annotate\_batch\_pre()}:

\begin{equation}
  \tilde{x}_{b,j}(t) = \frac{x_{b,j}(t) - \mu_{b,j}}{\sigma_{b,j} + \epsilon_{\mathrm{norm}}}, \quad j \in \{Z, N, E\},
  \label{eq:batch_trace_std_norm}
\end{equation}
where component mean $\mu_{b,j}$ and standard deviation $\sigma_{b,j}$ are computed
independently across the temporal dimension $t \in [1, L]$:
\begin{equation}
  \mu_{b,j} = \frac{1}{L} \sum_{t=1}^{L} x_{b,j}(t), \quad
  \sigma_{b,j} = \sqrt{\frac{1}{L} \sum_{t=1}^{L} \bigl( x_{b,j}(t) - \mu_{b,j} \bigr)^2},
\end{equation}
with $\epsilon_{\mathrm{norm}} = 10^{-6}$ preventing division by zero on dead or
disconnected channels. Performing standard normalization dynamically within the
GPU batch tensor graph guarantees exact numerical consistency with operational
real-time inference in \path{OGS/src/ogspicker.py}.

\subsection{Optimization Dynamics and Deterministic Seeding}
\label{subsec:optimization_dynamics}

Model parameters $\bm{\Theta}$ are updated using the Adam adaptive moment estimation
optimizer (\cite{KingmaBa2014}):
\begin{align}
  \mathbf{m}_t &= \beta_1 \mathbf{m}_{t-1} + (1 - \beta_1) \mathbf{g}_t, \\
  \mathbf{v}_t &= \beta_2 \mathbf{v}_{t-1} + (1 - \beta_2) \mathbf{g}_t^2, \\
  \hat{\mathbf{m}}_t &= \frac{\mathbf{m}_t}{1 - \beta_1^t}, \quad
  \hat{\mathbf{v}}_t = \frac{\mathbf{v}_t}{1 - \beta_2^t}, \\
  \bm{\Theta}_t &= \bm{\Theta}_{t-1} - \frac{\eta}{\sqrt{\hat{\mathbf{v}}_t} + \epsilon_{\mathrm{adam}}} \hat{\mathbf{m}}_t,
\end{align}
where $\mathbf{g}_t = \nabla_{\bm{\Theta}} \mathcal{L}$ denotes the stochastic gradient,
$\beta_1 = 0.9$ is the first moment decay, $\beta_2 = 0.999$ is the second moment
decay, $\epsilon_{\mathrm{adam}} = 10^{-8}$ is the smoothing term, and $\eta$ is
the learning rate.

The default configuration employs an initial learning rate $\eta = 10^{-2}$ for
broad parameter adaptation, adjustable via the CLI parameter \texttt{-lr / --learning\_rate}
down to $\eta = 10^{-3}$ for conservative fine-tuning of deep attention layers.
Mini-batches are assembled with $B = 256$ samples, providing a balanced trade-off
between gradient variance reduction and GPU memory occupancy.

To guarantee computational reproducibility across parallel DataLoader worker
subprocesses, worker initialization is tied to \texttt{seisbench.util.worker\_seeding}:
\begin{lstlisting}[language=python,basicstyle=\scriptsize\ttfamily]
train_loader = DataLoader(
    train_generator,
    batch_size=self.args.batch_size,
    shuffle=True,
    num_workers=self.args.workers,
    worker_init_fn=worker_seeding,
)
\end{lstlisting}
The \texttt{worker\_seeding} function derives a deterministic random seed for
each worker process by combining the master PyTorch random state with the worker
rank identifier, ensuring invariant augmentation sequences across experimental runs.

% =============================================================================
% SECTION 6: CHECKPOINTING, MODEL EVALUATION, AND VALIDATION STRATEGY
% =============================================================================
\section{Checkpointing, Model Evaluation, and Validation Strategy}
\label{sec:checkpointing_and_validation}

Robust training requires systematic evaluation against unseen validation data
to diagnose overfitting, guide early stopping, and preserve optimal model states.

\subsection{Dataset Partitioning: Train, Development, and Test Splits}
\label{subsec:dataset_partitioning}

The SeisBench \texttt{WaveformDataset.train\_dev\_test()} protocol partitions
the loaded dataset into three mutually disjoint subsets:
\begin{equation}
  \mathcal{D} = \mathcal{D}_{\mathrm{train}} \cup \mathcal{D}_{\mathrm{dev}} \cup \mathcal{D}_{\mathrm{test}},
  \quad \text{with} \quad \mathcal{D}_i \cap \mathcal{D}_j = \emptyset \; \forall i \neq j.
  \label{eq:dataset_partition}
\end{equation}

Partitioning is performed at the discrete earthquake event level rather than across
isolated traces. This spatial-temporal grouping prevents data leakage: all multi-station
recordings associated with a single seismic event are assigned entirely to one
split, preventing the neural network from memorizing source focal mechanisms or
event-specific wavefield signatures during training.

\subsection{Validation Protocol and Memory Optimization}
\label{subsec:validation_protocol}

At the conclusion of each training epoch, the validation method \texttt{\_validate()}
evaluates the model across the development generator $\mathcal{D}_{\mathrm{dev}}$:
\begin{lstlisting}[language=python,basicstyle=\scriptsize\ttfamily]
self.model.eval()
total_loss = 0.0
with torch.no_grad():
    for batch in dataloader:
        x = batch["X"].to(self.model.device)
        x_preproc = self.model.annotate_batch_pre(x, {})
        pred = self.model(x_preproc)
        total_loss += loss_fn(pred, batch["y"].to(self.model.device)).item()
return total_loss / max(num_batches, 1)
\end{lstlisting}

Two critical execution primitives govern this stage:
\begin{enumerate}
  \item \texttt{model.eval()}: Disables dropout layers and freezes batch
    normalization running statistics, ensuring deterministic evaluation.
  \item \texttt{torch.no_grad()}: Disables autograd gradient tape tracking,
    eliminating intermediate activation retention in GPU VRAM and speeding up
    validation inference by over $300\%$.
\end{enumerate}

\subsection{State-Preserving Checkpointing Architecture}
\label{subsec:checkpointing_arch}

The training engine implements a dual checkpoint retention strategy coordinated
by \texttt{\_save\_checkpoint()}:
\begin{enumerate}
  \item \textbf{Epoch-Wise State Dumps} (\path{checkpoint_epoch_<EEE>.pt}):
    At each epoch $e \in \{1, \dots, E\}$, an exhaustive checkpoint archive is
    persisted to the output directory:
    \begin{lstlisting}[language=python,basicstyle=\scriptsize\ttfamily]
checkpoint = {
    "epoch": epoch,
    "model_state_dict": self.model.state_dict(),
    "optimizer_state_dict": optimizer.state_dict(),
    "train_loss": train_loss,
    "val_loss": val_loss,
    "model_name": self.args.model,
    "dataset": self.args.dataset,
}
    \end{lstlisting}
    Retaining the complete optimizer state dictionary (\texttt{optimizer\_state\_dict})
    allows interrupted jobs on HPC clusters to resume without resetting Adam
    momentum and variance tracking buffers.
  \item \textbf{Optimal Model Tracking} (\path{best_model.pt}):
    The trainer continuously monitors the development set average loss
    $\bar{\mathcal{L}}_{\mathrm{val}}$. Whenever the current epoch yields a strictly
    lower loss than all preceding epochs ($\bar{\mathcal{L}}_{\mathrm{val}} < \mathcal{L}_{\mathrm{best}}$),
    the pure model state weights are updated:
    \begin{equation}
      \bm{\Theta}_{\mathrm{best}} \leftarrow \bm{\Theta}_e, \quad \mathcal{L}_{\mathrm{best}} \leftarrow \bar{\mathcal{L}}_{\mathrm{val}}.
      \label{eq:best_model_update}
    \end{equation}
    The resulting \path{best_model.pt} artifact represents the production-ready
    fine-tuned weights, directly loadable by \path{OGS/src/ogspicker.py} for
    inference across the operational pipeline.
\end{enumerate}

% =============================================================================
% SECTION 7: HPC WORKFLOW ORCHESTRATION ON CINECA LEONARDO
% =============================================================================
\section{HPC Workflow Orchestration on CINECA Leonardo}
\label{sec:hpc_orchestration}

Training and fine-tuning deep neural networks across continuous multi-year seismic
archives requires petascale computational throughput. The training workflow is
engineered for scalable execution on the \textbf{Leonardo} supercomputer hosted
at \ac{CINECA} under EuroHPC.

\subsection{Leonardo Booster Architecture and Hardware Acceleration}
\label{subsec:leonardo_booster}

The training pipeline is deployed on the Leonardo Booster partition, whose
architectural attributes were introduced in Chapter~\ref{chap:hpc}. Each Atos
BullSequana X2135 compute node integrates:
\begin{itemize}
  \item \textbf{Compute Accelerators}: Four NVIDIA Ampere A100-SXM4 GPUs with
    $64\,\unit{\giga\byte}$ of high-bandwidth HBM2e memory per GPU, delivering
    $1.55\,\unit{\tera\byte\per\second}$ of memory bandwidth and $19.5\,\unit{\tera\text{FLOPS}}$
    of single-precision (FP32) performance ($156\,\unit{\tera\text{FLOPS}}$ TF32).
  \item \textbf{Host Processor}: A 32-core Intel Xeon Platinum 8358 CPU
    ($2.60\,\unit{\giga\hertz}$, 64 threads) with $512\,\unit{\giga\byte}$ of
    DDR4-3200 host memory.
  \item \textbf{High-Speed Interconnect}: NVIDIA NVLink 3.0 inter-GPU fabric
    providing $600\,\unit{\giga\byte\per\second}$ bidirectional communication,
    complemented by dual 200G HDR InfiniBand adapters in Dragonfly+ topology.
  \item \textbf{Parallel Filesystem}: Shared high-speed Lustre scratch filesystem
    (\path{/leonardo_work/IscrC_AISeism/}) capable of delivering $> 100\,\unit{\giga\byte\per\second}$
    aggregate I/O throughput for HDF5 dataset ingestion.
\end{itemize}

\subsection{Parameterized Workflow Execution via GNU Make}
\label{subsec:makefile_trainer}

To streamline training campaigns and enforce operational reproducibility, the
training workflow is parameterized within the central orchestration Makefile
at \path{OGS/utils/Leonardo/Makefile}.

The \texttt{Trainer} target encapsulates model selection, baseline dataset,
hyperparameters, and directory paths:
\begin{lstlisting}[language=make,basicstyle=\scriptsize\ttfamily]
# Training defaults
EPOCHS        ?= 5
BATCH_SIZE    ?= 256
LEARNING_RATE ?= 1e-2
WORKERS       ?= 4

Trainer:
	@bash LAUNCHME.sh $(SLURM_TEMPLATE) 1 1 "Trainer_$(MODEL)_$(DATASET)" \
	  $(PYTHON_BIN) $(OGS_PATH)/src/ogstrainer.py \
	    -C $(CATALOG_PATH) \
	    -W $(WAVE_PATH) \
	    -D $(START_DATE) $(END_DATE) \
	    -m $(MODEL) \
	    -s $(DATASET) \
	    -b $(BATCH_SIZE) \
	    -e $(EPOCHS) \
	    -lr $(LEARNING_RATE) \
	    -w $(WORKERS) \
	    -o $(WORK_PATH)/models/$(MODEL)/$(DATASET)
\end{lstlisting}

Researchers can override any hyperparameter directly via CLI variables:
\begin{lstlisting}[language=bash,basicstyle=\scriptsize\ttfamily]
# Fine-tune PhaseNet on INSTANCE weights for 20 epochs with learning rate 1e-3
make Trainer MODEL=PhaseNet DATASET=INSTANCE EPOCHS=20 BATCH_SIZE=128 LEARNING_RATE=1e-3

# Fine-tune EQTransformer from STEAD global weights
make Trainer MODEL=EQTransformer DATASET=STEAD EPOCHS=15 BATCH_SIZE=64 LEARNING_RATE=5e-4
\end{lstlisting}

\subsection{SLURM Job Wrapping via LAUNCHME.sh}
\label{subsec:slurm_launchme}

Execution on Leonardo is managed through the dynamic SLURM job submission wrapper
\path{OGS/utils/Leonardo/LAUNCHME.sh}, interfacing with the job template
\path{OGS/utils/Leonardo/ktanakah.sh}.

When invoked with \texttt{bash LAUNCHME.sh <template> 1 1 <jobName> <command>},
the wrapper dynamically populates the SLURM allocation headers:
\begin{lstlisting}[language=bash,basicstyle=\scriptsize\ttfamily]
#SBATCH --job-name="Trainer_PhaseNet_INSTANCE"
#SBATCH --nodes=1
#SBATCH --tasks-per-node=1
#SBATCH --gres=gpu:1
#SBATCH --cpus-per-task=32
#SBATCH --account=IscrC_AISeism
#SBATCH --mem=490000MB
#SBATCH --partition=boost_usr_prod
#SBATCH --qos=boost_qos_dbg
#SBATCH --time=00:30:00
\end{lstlisting}
This allocates a dedicated NVIDIA A100 GPU (\texttt{--gres=gpu:1}) paired with
all 32 host CPU cores (\texttt{--cpus-per-task=32}) and $490\,\unit{\giga\byte}$
system RAM on the production Booster partition (\texttt{boost\_usr\_prod}).
Before launching Python, the wrapper invokes \path{ACTIVATEME.sh}, activating
the optimized Conda environment \texttt{SBC\_3.12} containing PyTorch compiled
with CUDA 12, cuDNN, and SeisBench libraries.

\subsection{End-to-End Pipeline DAG and Evidence Synthesis}
\label{subsec:pipeline_dag_evidence}

The complete training and domain adaptation workflow is synthesized in the
directed acyclic graph depicted in Figure~\ref{fig:training_dag}.

\begin{figure}[htbp]
  \centering
  \begin{tikzpicture}[
      >=Stealth,
      node distance=5mm and 6mm,
      every node/.style={font=\scriptsize},
      base/.style={rectangle, rounded corners=3pt, draw=black!80,
      line width=0.7pt, align=center, inner sep=4pt, minimum height=7mm},
      data/.style={base, fill=gray!10, draw=black!80, double, double
      distance=1pt, text width=26mm},
      proc/.style={base, fill=blue!8, draw=blue!80!black, text width=28mm},
      qc/.style={base, fill=teal!10, draw=teal!80!black, text width=26mm},
      planned/.style={base, fill=gray!4, draw=black!50, dash pattern=on 3.5pt off 2.5pt, text width=26mm, font=\scriptsize\itshape},
      arrow/.style={->, thick, draw=black!75},
      dashedarrow/.style={->, thick, dashed, draw=black!50}
    ]

    % Row 1: Data Extraction
    \node[data] (cat) {OGS Catalog\\(Weights $w \le 2$)};
    \node[data, right=of cat] (mseed) {Raw Continuous MSEED\\(\path{WAVE_PATH})};
    \node[proc, right=of mseed] (extract) {Stream Extraction\\($\pm 60\,\unit{\second}$ windowing)};

    % Row 2: Preprocessing & Serialization
    \node[proc, below=of extract] (clean) {Conditioning Pipeline\\(\texttt{get_clean_trace})};
    \node[data, below=of mseed] (hdf5) {SeisBench HDF5\\(\path{waveforms.hdf5})};
    \node[data, below=of cat] (meta) {Columnar Index\\(\path{metadata.csv})};

    % Row 3: Augmentation & Batching
    \node[proc, below=of hdf5] (gen) {SeisBench Generator\\(\texttt{WindowAroundSample})};
    \node[proc, right=of gen] (label) {Probabilistic Labeller\\($\sigma = 30$ samples)};
    \node[proc, right=of label] (batch) {DataLoader Workers\\($B=256$, seed init)};

    % Row 4: GPU Acceleration & Training
    \node[proc, below=of batch] (norm) {GPU Normalization\\(\texttt{annotate_batch_pre})};
    \node[proc, below=of label] (model) {Neural Forward Pass\\(PhaseNet / EQT)};
    \node[proc, below=of gen] (loss) {Vector Cross-Entropy\\(\texttt{loss_fn}, $\epsilon=10^{-5}$)};

    % Row 5: Checkpointing & Validation
    \node[proc, below=of loss] (opt) {Adam Optimization\\($\eta \in [10^{-3}, 10^{-2}]$)};
    \node[data, below=of model] (ckpt) {Checkpoint Archives\\(\path{best_model.pt})};
    \node[planned, below=of norm] (bench) {[Proposed] Regional\\Evaluation Suite};

    % Connections
    \draw[arrow] (cat) -- (extract);
    \draw[arrow] (mseed) -- (extract);
    \draw[arrow] (extract) -- (clean);
    \draw[arrow] (clean) -- (hdf5);
    \draw[arrow] (clean) -- (meta);
    \draw[arrow] (hdf5) -- (gen);
    \draw[arrow] (meta) -- (gen);
    \draw[arrow] (gen) -- (label);
    \draw[arrow] (label) -- (batch);
    \draw[arrow] (batch) -- (norm);
    \draw[arrow] (norm) -- (model);
    \draw[arrow] (model) -- (loss);
    \draw[arrow] (loss) -- (opt);
    \draw[arrow] (opt) -- (ckpt);
    \draw[dashedarrow] (ckpt) -- (bench);

    % Subsystem Boxes
    \begin{pgfonlayer}{background}
      \node[draw=blue!40, fill=blue!2, dashed, rounded corners=5pt,
      fit=(clean) (extract), inner sep=4pt] (ingestbox) {};
      \node[anchor=north west, font=\tiny\bfseries\color{blue!70!black}] at
      (ingestbox.north west) {Stream Conditioning};

      \node[draw=teal!40, fill=teal!2, dashed, rounded corners=5pt,
      fit=(norm) (model) (loss) (opt), inner sep=4pt] (trainbox) {};
      \node[anchor=south west, font=\tiny\bfseries\color{teal!70!black}] at
      (trainbox.south west) {Leonardo Booster GPU Execution};
    \end{pgfonlayer}

  \end{tikzpicture}
  \caption{Directed Acyclic Graph (DAG) of the regional deep-learning training
    and domain adaptation pipeline. Double-bordered gray boxes denote immutable
    persistent filesystem storage (catalog records, miniSEED streams, HDF5 containers,
    and checkpoint binaries); solid blue boxes indicate demonstrated software modules
    implemented in \path{OGS/src/ogstrainer.py}; dashed gray boxes indicate planned
    validation benchmarks. Notice that stream conditioning and HDF5 serialization
    occur prior to stochastic augmentation, ensuring zero-copy GPU batch loading.}
  \label{fig:training_dag}
\end{figure}

Table~\ref{tab:training_hyperparameters} provides an exhaustive specification
of all operational hyperparameters, network dimensions, and optimization parameters
governing PhaseNet and EQTransformer training within the Leonardo Booster environment.

\begin{table}[htbp]
  \centering
  \small
  \caption{Operational hyperparameters, architectural dimensions, and optimization
    configurations for regional domain adaptation on Leonardo Booster.}
  \label{tab:training_hyperparameters}
  \begin{threeparttable}
    \begin{tabularx}{\linewidth}{l c c X}
      \toprule
      \textbf{Configuration Parameter} & \textbf{PhaseNet Setup} & \textbf{EQTransformer Setup} & \textbf{Operational Rationale \& Algorithmic Constraint} \\
      \midrule
      Input Sequence Length ($L$) & $3{,}001$ samples & $6{,}000$ samples & $30.01\,\unit{\second}$ vs $60.00\,\unit{\second}$ at $100\,\unit{\hertz}$; captures full $S$-coda. \\
      Input Waveform Channels ($C$)& $3$ ($Z, N, E$)   & $3$ ($Z, N, E$)   & Normalized 3-component particle motion velocities. \\
      Pretrained Weight Baselines & \texttt{instance} & \texttt{stead}     & Warm-start representations from Italian or global corpora. \\
      Mini-Batch Size ($B$)       & $256$             & $128$ / $256$      & Balances GPU VRAM saturation with gradient variance reduction. \\
      Optimizer Class             & Adam              & Adam               & First moment $\beta_1=0.9$, second moment $\beta_2=0.999$. \\
      Base Learning Rate ($\eta$) & $1.0 \times 10^{-2}$& $1.0 \times 10^{-2}$& Standard CLI default; fine-tunable down to $10^{-3}$. \\
      Loss Regularization ($\epsilon$) & $1.0 \times 10^{-5}$ & $1.0 \times 10^{-5}$ & Numerical stability constant in vector cross-entropy. \\
      Label Target Standard Dev ($\sigma$) & $30$ samples & $30$ samples & $0.30\,\unit{\second}$ Gaussian envelope matching pick quality $w \le 2$. \\
      Preprocessing Bandpass Filter & $0.1\text{--}20.0\,\unit{\hertz}$ & $0.1\text{--}20.0\,\unit{\hertz}$ & 4th-order zero-phase Butterworth; isolates seismic signal. \\
      Pick Quality Weight Threshold & $w \le 2$         & $w \le 2$          & Discards emergent picks with timing uncertainty $> 0.2\,\unit{\second}$. \\
      DataLoader Workers ($N_w$)  & $4$               & $4$                & Asynchronous multi-process loading with \texttt{worker_seeding}. \\
      Hardware Allocation (SLURM) & $1 \times$ A100 GPU & $1 \times$ A100 GPU & $64\,\unit{\giga\byte}$ HBM2e, 32 CPU cores, $490\,\unit{\giga\byte}$ RAM on Booster. \\
      \bottomrule
    \end{tabularx}
    \begin{tablenotes}[flushleft]
      \footnotesize
      \item Target \& Scope: Algorithmic parameters for neural network fine-tuning
        configured in \path{OGS/src/ogstrainer.py} and submitted via \path{OGS/utils/Leonardo/Makefile}.
      \item Notice the difference in input sequence length ($3{,}001$ vs $6{,}000$ samples),
        reflecting EQTransformer's recurrent attention mechanism over extended coda sequences.
      \item Evidence Anchor: Formulated in \texttt{ogstrainer.py} functions \texttt{\_build\_augmentations()},
        \texttt{loss\_fn()}, and Makefile target \texttt{Trainer}.
    \end{tablenotes}
  \end{threeparttable}
\end{table}

\section{Summary and Outlook}
\label{sec:training_summary}

This chapter has established the mathematical principles, software architecture,
and HPC execution protocol for deep-learning seismic phase picker fine-tuning and
regional domain adaptation. By integrating analyst-curated arrival picks from the
\ac{OGS} bulletin with continuous multi-channel waveform streams, the system
enforces strict pick uncertainty constraints ($w \le 2$, $\sigma_t \le 0.2\,\unit{\second}$),
deterministic signal conditioning, stochastic multi-window augmentation, and
continuous Gaussian label smoothing.

The derived vector cross-entropy loss function guarantees numerically stable
gradient updates on GPU accelerators, while the dual checkpoint retention protocol
safeguards optimal model states for operational catalog generation. Deploying this
framework through parameterized Makefile rules and SLURM wrappers on the CINECA
Leonardo Booster supercomputer provides the computational foundation for
evaluating regional domain adaptation across the complex tectonic boundaries of
Northeastern Italy, paving the way for high-resolution automated seismicity
monitoring.

